% !TEX root = ../main.tex

\documentclass[../main.tex]{subfiles}

\graphicspath{{\subfix{../images/}}}

\begin{document}

\subsection{Функция \texttt{integrate\_imu()}}

Весь навигационный алгоритм сосредоточен в функции \texttt{integrate\_imu()}
файла \texttt{imu\_visualizer.cpp}. Она вызывается из потока визуализации для
каждого пакета, извлечённого из очереди. Шаг интегрирования $\Delta t$
вычисляется как разность временны'х меток двух последовательных пакетов.
 
\textbf{Блок 1.} Калибровка сырых отсчётов по формуле~\eqref{eq:calib}:
 
\begin{lstlisting}[caption={Калибровка сырых показаний}]
double ax = (raw.xacc  - Calib::AX_BIAS) * Calib::ACCEL_SCALE;
double ay = (raw.yacc  - Calib::AY_BIAS) * Calib::ACCEL_SCALE;
double az = (raw.zacc  - Calib::AZ_BIAS) * Calib::ACCEL_SCALE;
double gx = (raw.xgyro - Calib::GX_BIAS) * Calib::GYRO_SCALE;
double gy = (raw.ygyro - Calib::GY_BIAS) * Calib::GYRO_SCALE;
double gz = (raw.zgyro - Calib::GZ_BIAS) * Calib::GYRO_SCALE;
\end{lstlisting}
 
\textbf{Блок 2.} Кватернионное интегрирование через ось-угол по
формулам~\eqref{eq:dq}--\eqref{eq:qupdate}:
 
\begin{lstlisting}[caption={Кватернионное интегрирование гироскопа}]
double theta = std::sqrt(gx*gx + gy*gy + gz*gz) * dt;
if (theta > 1e-10) {
    double norm_w = std::sqrt(gx*gx + gy*gy + gz*gz);
    double s   = std::sin(theta * 0.5) / norm_w;
    double c   = std::cos(theta * 0.5);
    double dqw = c,    dqx = gx*s, dqy = gy*s, dqz = gz*s;
 
    // q_new = q (x) dq
    double qw = n.qw*dqw - n.qx*dqx - n.qy*dqy - n.qz*dqz;
    double qx = n.qw*dqx + n.qx*dqw + n.qy*dqz - n.qz*dqy;
    double qy = n.qw*dqy - n.qx*dqz + n.qy*dqw + n.qz*dqx;
    double qz = n.qw*dqz + n.qx*dqy - n.qy*dqx + n.qz*dqw;
 
    double norm = std::sqrt(qw*qw + qx*qx + qy*qy + qz*qz);
    n.qw=qw/norm; n.qx=qx/norm; n.qy=qy/norm; n.qz=qz/norm;
}
\end{lstlisting}
 
\textbf{Блок 3.} Перевод ускорения в мировую систему координат,
вычитание гравитации и применение ZUPT:
 
\begin{lstlisting}[caption={Мировое ускорение, компенсация гравитации и ZUPT}]
// Матрица поворота (строки r00..r22) из кватерниона
double ax_w = r00*ax + r01*ay + r02*az;
double ay_w = r10*ax + r11*ay + r12*az;
double az_w = r20*ax + r21*ay + r22*az;
az_w -= (-9.81);         // вычитаем гравитацию
 
// ZUPT: гироскоп в покое -> мягко гасим скорость
double gyro_norm = std::sqrt(gx*gx + gy*gy + gz*gz);
if (gyro_norm < Calib::ZUPT_GYRO_THRESH) {
    n.vx *= 0.85; n.vy *= 0.85; n.vz *= 0.85;
    if (std::abs(ax_w) < 0.3) ax_w = 0;
    if (std::abs(ay_w) < 0.3) ay_w = 0;
    if (std::abs(az_w) < 0.3) az_w = 0;
}
// Двойное интегрирование
n.vx += ax_w*dt; n.vy += ay_w*dt; n.vz += az_w*dt;
n.px += n.vx*dt; n.py += n.vy*dt; n.pz += n.vz*dt;
\end{lstlisting}
 
\subsection{Вычисление углов Эйлера}
 
Завершает функцию вычисление углов Эйлера по
формулам~\eqref{eq:roll}--\eqref{eq:yaw}. Аргумент \texttt{asin}
зажимается в интервал $[-1,1]$ для защиты от накопленной численной
погрешности:
 
\begin{lstlisting}[caption={Углы Эйлера из кватерниона (ZYX)}]
const double R2D = 180.0 / M_PI;
n.roll  = std::atan2(2*(qw*qx + qy*qz),
                     1 - 2*(qx*qx + qy*qy)) * R2D;
double sp = std::max(-1.0, std::min(1.0, 2*(qw*qy - qz*qx)));
n.pitch = std::asin(sp) * R2D;
n.yaw   = std::atan2(2*(qw*qz + qx*qy),
                     1 - 2*(qy*qy + qz*qz)) * R2D;
\end{lstlisting}
	
\end{document}
